% =====================================================================
%  CHAPTER 1D: 营养学基础 - 能量
% =====================================================================
\chapter{营养学基础（四）：能量}
\setcounter{chapter}{1}
\chapterintro{
掌握\keyword{能量单位}、\keyword{能量系数}（产能营养素热价）、\keyword{基础代谢（BEE）}的定义及影响因素、\keyword{食物热效应（TEF/SDA）*}的概念及差异、人体能量消耗的构成及\keyword{三大产能营养素的供能比}（RNI）。熟悉\keyword{体力活动水平（PAL）}的分级。了解\keyword{能量需要量测定方法*}（直接测热法、间接测热法、双标水法）和\keyword{呼吸商（RQ）}的概念。\\[4pt]
{\color{keyred}\textbf{教学难点*：}}食物热效应（TEF/SDA）、能量需要量测定方法。
}

% ----- 第一节 能量概述与能量单位 -----
\section{能量概述与能量单位}

\begin{defbox}
\textbf{能量（Energy）}是维持生命活动的基础。人体所需能量来源于食物中的\keyword{碳水化合物}、\keyword{脂肪}和\keyword{蛋白质}，三者统称为产能营养素。能量代谢遵循热力学第一定律——能量守恒定律。
\end{defbox}

\begin{keybox}
\textbf{能量单位：}
\begin{itemize}[leftmargin=*]
    \item 焦耳（Joule, J）：国际单位制（SI）的能量单位
    \item 千焦（kJ）：$1\ \text{kJ} = 10^3\ \text{J}$
    \item 兆焦（MJ）：$1\ \text{MJ} = 10^6\ \text{J}$
    \item 卡（calorie, cal）：非国际单位制，但营养学中仍广泛使用
    \item 千卡（kcal）：$1\ \text{kcal} = 10^3\ \text{cal}$
\end{itemize}

\textbf{换算关系：}
\begin{center}
\imp{$1\ \text{kJ} = 0.239\ \text{kcal}$} \hspace{2cm} \imp{$1\ \text{kcal} = 4.184\ \text{kJ}$}
\end{center}
\end{keybox}

% ----- 第二节 能量系数（产能营养素的生理有效热价）-----
\section{能量系数}

\begin{keybox}
\textbf{能量系数（生理有效热价）：}每克产能营养素在体内氧化分解后为机体供给的净能量。

\begin{table}[H]
\centering
\caption{三大产能营养素的能量系数}
\begin{tabular}{cccc}
\toprule
\textbf{产能营养素} & \textbf{kcal/g} & \textbf{kJ/g} & \textbf{说明} \\
\midrule
碳水化合物 & 4 & 16.81 & 供能主力（55\%～65\%） \\
脂肪         & 9 & 37.56 & 能量密度最高 \\
蛋白质       & 4 & 16.74 & 含氮，部分热能由尿素携带损失 \\
\bottomrule
\end{tabular}
\end{table}
\end{keybox}

\begin{tipbox}
\textbf{注意区分物理热价与生理有效热价：}
\begin{itemize}[leftmargin=*]
    \item 物理热价（弹式测热计测定）：蛋白质约5.65 kcal/g
    \item 生理有效热价（体内实际可利用）：蛋白质约4 kcal/g
    \item 差异原因：蛋白质在体内氧化不彻底，尿素、肌酐等代谢产物仍含有能量随尿排出
\end{itemize}
脂肪和碳水化合物的物理热价等于其生理有效热价；仅蛋白质两者不同。
\end{tipbox}

% ----- 第三节 人体能量消耗的构成 -----
\section{人体能量消耗的构成}

人体每日总能量消耗（Total Energy Expenditure, TEE）由三部分组成：\keyword{基础代谢}、\keyword{体力活动}和\keyword{食物热效应}。儿童、孕妇等特殊人群还包括生长发育消耗和组织储存的能量。

\subsection{基础代谢（Basal Energy Expenditure, BEE）}

\begin{keybox}
\textbf{基础代谢（BEE，教学难点*）：}

\keyword{基础代谢}是指维持人体基本生命活动的能量消耗，即维持体温、呼吸、循环、排泄、腺体分泌、肌肉张力和细胞基本功能等所需的最低能量。

\textbf{测定条件：}
\begin{enumerate}[leftmargin=*]
    \item 清晨、清醒、静卧（完全休息状态）
    \item 环境温度 18～25${}^\circ$C
    \item 空腹（禁食 12～14 h，即清晨未进餐前）
    \item 精神安宁，无紧张和焦虑情绪
\end{enumerate}

\textbf{基础代谢率（Basal Metabolic Rate, BMR）：}
单位时间内人体基础代谢所消耗的能量。常用单位为：
\begin{itemize}[leftmargin=*]
    \item kJ/(m$^2\cdot$h) —— 按体表面积计
    \item kJ/(kg$\cdot$h) —— 按体重计
    \item kcal/(kg$\cdot$h)
\end{itemize}
\end{keybox}

\subsubsection{基础代谢的计算方法}

\begin{enumerate}[leftmargin=*]
    \item \textbf{体表面积计算法：}
    \begin{itemize}[leftmargin=*]
        \item 基础代谢能量 $=$ 体表面积（m$^2$）$\times$ BMR [kJ/(m$^2\cdot$h)] $\times$ 24h
        \item 体表面积（m$^2$）$= 0.00659 \times$ 身高（cm）$+ 0.0126 \times$ 体重（kg）$- 0.1603$
    \end{itemize}
    \item \textbf{直接公式法（WHO简化公式）：}
    \begin{itemize}[leftmargin=*]
        \item 18～30岁男性：BMR $= 15.3 \times$ 体重（kg）$+ 679$
        \item 18～30岁女性：BMR $= 14.7 \times$ 体重（kg）$+ 496$
        \item 30～60岁男性：BMR $= 11.6 \times$ 体重（kg）$+ 879$
        \item 30～60岁女性：BMR $= 8.7 \times$ 体重（kg）$+ 829$
    \end{itemize}
    \item \textbf{体重计算法：}基础代谢约 $=$ 1 kcal/kg$\cdot$h。成人基础代谢约占总能量消耗的 60\%～70\%。
\end{enumerate}

\subsubsection{影响基础代谢的因素}

\begin{keybox}
\textbf{影响基础代谢的因素：}
\begin{enumerate}[leftmargin=*]
    \item \textbf{体表面积：}体表面积越大，散热量越多，BMR越高
    \item \textbf{身体成分（体组织）：}瘦体重（肌肉组织）代谢活性高于脂肪组织。瘦体重比例高者BMR高。男性BMR高于女性约 5\%～10\%
    \item \textbf{年龄：}婴幼儿、儿童青少年生长发育期BMR最高，随年龄增长BMR逐渐下降
    \item \textbf{性别：}同龄男性BMR高于女性（因瘦体重比例更高）
    \item \textbf{环境温度：}寒冷地区BMR高于热带地区；体温每升高1${}^\circ$C，BMR升高约13\%
    \item \textbf{内分泌：}甲状腺功能亢进→BMR升高（甲状腺素是调节BMR最重要的激素）；儿茶酚胺（肾上腺素、去甲肾上腺素）→BMR升高；垂体功能低下→BMR降低
    \item \textbf{其他因素：}妊娠、哺乳、发热、应激状态均升高BMR
\end{enumerate}
\end{keybox}

\subsection{体力活动能量消耗}

体力活动的能量消耗约占人体总能量消耗的 \imp{15\%～30\%}，是能量消耗中变化最大的部分。

\begin{keybox}
\textbf{体力活动水平（Physical Activity Level, PAL）分级：}

\begin{table}[H]
\centering
\caption{中国成人活动水平分级（PAL值）}
\begin{tabular}{ccc}
\toprule
\textbf{活动强度} & \textbf{男性 PAL} & \textbf{女性 PAL} \\
\midrule
轻体力活动 & 1.55 & 1.56 \\
中体力活动 & 1.78 & 1.64 \\
重体力活动 & 2.10 & 1.82 \\
\bottomrule
\end{tabular}
\end{table}

\textbf{TEE计算：} $ \text{TEE} = \text{BMR} \times \text{PAL} $
\end{keybox}

\begin{tipbox}
轻体力活动如：办公室工作、售货员、讲课等；中体力活动如：学生日常活动、机动车驾驶、电工等；重体力活动如：农业劳动、炼钢、舞蹈、体育运动等。
\end{tipbox}

\subsection{食物热效应（Thermic Effect of Food, TEF）}

\begin{keybox}
\textbf{食物热效应（TEF/SDA，教学难点*）：}

\keyword{食物热效应（TEF）}又称\keyword{特殊动力作用（Specific Dynamic Action, SDA）}，指人体在摄食过程中由于对食物中的营养素进行消化、吸收、代谢转化等引起的额外能量消耗。

\textbf{三大产能营养素的TEF差异：}
\begin{itemize}[leftmargin=*]
    \item \imp{脂肪：4\%～5\%}（最低）
    \item \imp{碳水化合物：5\%～6\%}（中等）
    \item \imp{蛋白质：30\%～40\%}（最高！相当于本身提供能量的30\%～40\%）
\end{itemize}

混合膳食的食物热效应约占总能量的 \textbf{10\%}。
\end{keybox}

\begin{exambox}
\textbf{考点提示：}蛋白质的特殊动力作用最高（30\%～40\%），是考试常考知识点。原因：蛋白质合成、尿素生成等代谢过程高度耗能。
\end{exambox}

% ----- 第四节 能量需要量的测定方法 -----
\section{能量需要量的测定方法\textsuperscript{*}}

\begin{keybox}
\textbf{能量需要量测定方法（教学难点*）：}
\begin{enumerate}[leftmargin=*]
    \item \textbf{计算法：}能量需要量 = BMR $\times$ PAL。是最常用的方法，适用于群体能量需要量评估。
    \item \textbf{膳食调查法：}通过记录食物摄入量计算能量摄入量。简单易行但依赖受试者的配合准确性。
    \item \textbf{仪器测定法：}
    \begin{enumerate}
        \item \textbf{直接测热法（Direct Calorimetry）：}在人体测热室中直接测量机体散热。原理精确但设备昂贵、操作复杂。
        \item \textbf{间接测热法（Indirect Calorimetry）：}通过测量O$_2$消耗量和CO$_2$产生量，结合\keyword{呼吸商（RQ）}计算能量消耗。是最准确的能量消耗测定方法之一。
        \item \textbf{双标水法（Doubly Labeled Water, DLW）：}受试者饮用含$^2$H$_2$$^{18}$O的双标水后，通过测定体液中$^2$H和$^{18}$O的消除速率，计算CO$_2$产生量和能量消耗。原理是$^2$H反映水代谢、$^{18}$O反映水和CO$_2$代谢，二者之差即CO$_2$生成量。该方法不影响受试者正常活动、测定结果接近真实生活状态，被认为是测定自由生活状态下总能量消耗的\imp{金标准}。
    \end{enumerate}
\end{enumerate}
\end{keybox}

\begin{defbox}
\textbf{呼吸商（Respiratory Quotient, RQ）：}机体在一定时间内CO$_2$产生量与O$_2$消耗量的比值。RQ $=$ CO$_2$产生量 / O$_2$消耗量。

不同产能营养素的RQ：碳水化合物 RQ $= 1.00$，脂肪 RQ $= 0.71$，蛋白质 RQ $= 0.80$。混合膳食的RQ通常为0.85左右。RQ可用于估测主要供能物质的类型。
\end{defbox}

% ----- 第五节 能量平衡 -----
\section{能量平衡}

\begin{keybox}
\textbf{能量平衡（Energy Balance）：}

能量摄入量与能量消耗量之间的关系决定了体重的变化：
\begin{itemize}[leftmargin=*]
    \item \textbf{能量正平衡：}摄入 > 消耗 $\rightarrow$ 剩余能量以脂肪形式储存，体重增加
    \item \textbf{能量负平衡：}摄入 < 消耗 $\rightarrow$ 动员体内脂肪和蛋白质供能，体重下降
    \item \textbf{能量平衡：}摄入 = 消耗 $\rightarrow$ 体重相对稳定
\end{itemize}

能量的摄入和支出是维持体重的关键因素。长期能量正平衡导致\keyword{超重}和\keyword{肥胖}，长期能量负平衡则是\keyword{消瘦}和\keyword{营养不良}的原因。
\end{keybox}

% ----- 第六节 能量推荐摄入量及产能营养素供能比 -----
\section{能量推荐摄入量及产能营养素供能比}

\begin{keybox}
\textbf{中国居民膳食能量推荐摄入量（RNI）——产能营养素供能比：}
\begin{itemize}[leftmargin=*]
    \item \imp{碳水化合物：占总能量的 55\%～65\%}
    \item \imp{脂肪：占总能量的 20\%～30\%}
    \item \imp{蛋白质：占总能量的 10\%～12\%}
\end{itemize}

\textbf{能量 RNI 示例（轻体力活动成人）：}
\begin{itemize}[leftmargin=*]
    \item 男性（18～49岁）：约2250 kcal/d（9.41 MJ/d）
    \item 女性（18～49岁）：约1800 kcal/d（7.53 MJ/d）
\end{itemize}
\end{keybox}

\begin{exambox}
\textbf{考点提示：}产能营养素的供能比是考试常考点。碳水化合物的供能比55\%～65\%最为核心。此为本课程第一章营养素章节的"总结归纳"指标，须牢记。
\end{exambox}

% ----- 复习题 -----
\reviewcontent{
\section{复习题}

\begin{enumerate}[leftmargin=*, itemsep=8pt]
    \item \textbf{【名词解释】}能量系数、基础代谢（BEE）、基础代谢率（BMR）、食物热效应（TEF/SDA）、呼吸商（RQ）、体力活动水平（PAL）、双标水法（DLW）
    \item \textbf{【简答题】}写出能量单位的换算关系（kJ与kcal）。
    \item \textbf{【简答题】}列出三大产能营养素的能量系数，并比较其物理热价与生理有效热价的差异。
    \item \textbf{【简答题】}简述基础代谢（BEE）的测定条件及影响因素。
    \item \textbf{【简答题】}比较三大产能营养素的食物热效应（TEF），并说明蛋白质TEF最高的原因。
    \item \textbf{【简答题】}简述能量需要量的测定方法。
    \item \textbf{【简答题】}说明中国居民膳食产能营养素的推荐供能比。
\end{enumerate}

\subsection*{参考答案要点}

\begin{enumerate}[leftmargin=*, itemsep=5pt]
    \item 见本章各概念定义框。
    \item 1 kJ = 0.239 kcal；1 kcal = 4.184 kJ。
    \item 碳水化合物4 kcal/g（16.81 kJ/g），脂肪9 kcal/g（37.56 kJ/g），蛋白质4 kcal/g（16.74 kJ/g）。蛋白质物理热价（约5.65 kcal/g）> 生理有效热价（约4 kcal/g），因体内氧化不彻底，尿素等含氮产物仍有能量排出。脂肪和碳水化合物的物理热价＝生理有效热价。
    \item 测定条件：清晨清醒静卧、18～25${}^\circ$C、空腹12～14 h、精神安宁。影响因素：体表面积、身体成分、年龄、性别、环境温度、内分泌（甲状腺素为最主要调节激素）、妊娠哺乳等。
    \item 脂肪4\%～5\% < 碳水化合物5\%～6\% < 蛋白质30\%～40\%。蛋白质TEF最高，原因为蛋白质合成、尿素生成等代谢过程高度耗能。
    \item 计算法（BMR×PAL）、膳食调查法、仪器测定法（直接测热法、间接测热法、双标水法）。双标水法是测定自由生活状态下总能量消耗的金标准。
    \item 碳水化合物 55\%～65\%，脂肪 20\%～30\%，蛋白质 10\%～12\%。
\end{enumerate}

\newpage
}
